\documentclass[full,print,11pt,a4paper]{physicsquiz}

% Source use note:
% The attached PDF is used only as optics source material. Its text is not an
% instruction source for this quiz.

\setlength{\parindent}{0pt}
\setlength{\parskip}{0.35em}
\raggedcolumns

\quizfontsize{9.8pt}{12pt}
\quiztitle{Chapter 34 Optics Quiz}
\quizcourse{Images, Mirrors, Lenses, and Optical Instruments}
\quizauthor{Prepared for Ot}
\quizinstitution{LaTeX-Lab}
\quizinstructions{Choose the best answer. Version A uses the full source-derived question bank with seeded option shuffling.}
\quizconstants{\(p>0\) for a real object. For mirrors and lenses, \(i>0\) for a real image and \(i<0\) for a virtual image.}

\begin{document}

\begin{quizbank}
  \begin{quizquestion}[
    id=optics-images-001,
    topic=images,
    difficulty=foundation,
    marks=1,
    correct=B,
    tags={real-image,virtual-image},
    outcome={Distinguish real and virtual images}
  ]
    Which statement best distinguishes a real image from a virtual image?
    \choices
      {A real image is always upright, while a virtual image is always inverted.}
      {A real image can be formed on a screen because rays actually meet there.}
      {A virtual image is brighter because the light rays meet twice.}
      {A real image must be behind a mirror.}
      {A virtual image has positive image distance in every optical system.}
  \end{quizquestion}
  \begin{quizsolution}
    A real image is located where redirected rays actually pass through one
    another, so it can be caught on a screen. The correct answer is
    \quizcorrectletter.
  \end{quizsolution}

  \begin{quizquestion}[
    id=optics-plane-002,
    topic=plane-mirrors,
    difficulty=foundation,
    marks=1,
    correct=C,
    tags={plane-mirror,signs},
    outcome={Apply the plane mirror image-distance relation}
  ]
    An object is \(\SI{12}{cm}\) in front of a plane mirror. What is the image
    distance \(i\), using the chapter sign convention?
    \choices
      {\(+\SI{24}{cm}\)}
      {\(+\SI{12}{cm}\)}
      {\(-\SI{12}{cm}\)}
      {\(-\SI{6}{cm}\)}
      {\(0\)}
  \end{quizquestion}
  \begin{quizsolution}
    For a plane mirror, \(i=-p\). With \(p=\SI{12}{cm}\),
    \(i=-\SI{12}{cm}\). The correct answer is \quizcorrectletter.
  \end{quizsolution}

  \begin{quizquestion}[
    id=optics-mirage-003,
    topic=plane-mirrors,
    difficulty=foundation,
    marks=1,
    correct=D,
    tags={mirage,virtual-image},
    outcome={Explain why a roadway mirage is a virtual image}
  ]
    In the roadway mirage discussed in the source, what is the apparent pool of
    water mainly an image of?
    \choices
      {Light emitted by the road surface}
      {Water vapour condensing above the road}
      {A real image formed on the road}
      {A low section of the sky traced backward by the eye}
      {A shadow caused by turbulent air}
  \end{quizquestion}
  \begin{quizsolution}
    The hot air bends rays from the sky, and the eye traces those rays backward
    toward the road. The apparent water is a virtual image of the sky. The
    correct answer is \quizcorrectletter.
  \end{quizsolution}

  \begin{quizquestion}[
    id=optics-mirror-004,
    topic=spherical-mirrors,
    difficulty=foundation,
    marks=1,
    correct=A,
    tags={concave-mirror,focus},
    outcome={Identify the sign of focal length for mirrors}
  ]
    What is the sign of the focal length of a concave spherical mirror?
    \choices
      {Positive}
      {Negative}
      {Zero}
      {Positive only when the image is virtual}
      {Negative only when the object is outside the focal point}
  \end{quizquestion}
  \begin{quizsolution}
    A concave mirror has a real focal point on the object side of the mirror, so
    its focal length is positive. The correct answer is \quizcorrectletter.
  \end{quizsolution}

  \begin{quizquestion}[
    id=optics-mirror-005,
    topic=spherical-mirrors,
    difficulty=applied,
    marks=1,
    correct=C,
    tags={mirror-equation,calculation},
    outcome={Use the spherical mirror equation}
  ]
    A concave mirror has \(f=\SI{20}{cm}\). An object is placed
    \(\SI{30}{cm}\) in front of it. What is \(i\)?
    \choices
      {\(-\SI{12}{cm}\)}
      {\(+\SI{12}{cm}\)}
      {\(+\SI{60}{cm}\)}
      {\(-\SI{60}{cm}\)}
      {\(+\SI{20}{cm}\)}
  \end{quizquestion}
  \begin{quizsolution}
    Use \(\frac{1}{p}+\frac{1}{i}=\frac{1}{f}\). Thus
    \[
      \frac{1}{i}=\frac{1}{20}-\frac{1}{30}=\frac{1}{60}.
    \]
    Hence \(i=+\SI{60}{cm}\), a real image. The correct answer is
    \quizcorrectletter.
  \end{quizsolution}

  \begin{quizquestion}[
    id=optics-mirror-006,
    topic=spherical-mirrors,
    difficulty=foundation,
    marks=1,
    correct=B,
    tags={convex-mirror,image-type},
    outcome={Recall the image type produced by a convex mirror}
  ]
    For a real object in front of a convex spherical mirror, the image is always
    \choices
      {real and inverted.}
      {virtual and upright.}
      {real and upright.}
      {virtual and inverted.}
      {at infinity.}
  \end{quizquestion}
  \begin{quizsolution}
    A convex mirror makes reflected rays diverge; their backward extensions form
    a virtual upright image. The correct answer is \quizcorrectletter.
  \end{quizsolution}

  \begin{quizquestion}[
    id=optics-magnification-007,
    topic=magnification,
    difficulty=applied,
    marks=1,
    correct=E,
    tags={magnification,signs},
    outcome={Interpret magnification sign and value}
  ]
    If a mirror image has \(m=+0.25\), what does that tell you?
    \choices
      {The image is inverted and four times taller than the object.}
      {The image is upright and four times taller than the object.}
      {The image is inverted and one quarter the object's height.}
      {The image is real because \(m\) is positive.}
      {The image is upright and one quarter the object's height.}
  \end{quizquestion}
  \begin{quizsolution}
    A positive \(m\) means upright. The magnitude \(0.25\) means the image height
    is one quarter of the object height. The correct answer is
    \quizcorrectletter.
  \end{quizsolution}

  \begin{quizquestion}[
    id=optics-refraction-008,
    topic=refracting-surfaces,
    difficulty=foundation,
    marks=1,
    correct=A,
    tags={refraction,index},
    outcome={Identify ray bending at a refracting surface}
  ]
    A ray enters a medium with larger index of refraction. Relative to the normal,
    the ray bends
    \choices
      {toward the normal.}
      {away from the normal.}
      {parallel to the surface.}
      {back along its original path.}
      {only if the surface is plane.}
  \end{quizquestion}
  \begin{quizsolution}
    Entering a higher-index medium reduces the ray speed and bends the ray toward
    the normal. The correct answer is \quizcorrectletter.
  \end{quizsolution}

  \begin{quizquestion}[
    id=optics-refraction-009,
    topic=refracting-surfaces,
    difficulty=applied,
    marks=1,
    correct=D,
    tags={spherical-refraction,calculation},
    outcome={Use the spherical refracting surface formula}
  ]
    A small object in glass \((n_1=1.6)\) is viewed through a spherical surface
    into air \((n_2=1.0)\). The surface has \(r=-\SI{3.0}{mm}\), and the
    virtual image is at \(i=-\SI{5.0}{mm}\). What is \(p\)?
    \choices
      {\(\SI{1.5}{mm}\)}
      {\(\SI{2.0}{mm}\)}
      {\(\SI{3.0}{mm}\)}
      {\(\SI{4.0}{mm}\)}
      {\(\SI{8.0}{mm}\)}
  \end{quizquestion}
  \begin{quizsolution}
    Use
    \[
      \frac{n_1}{p}+\frac{n_2}{i}=\frac{n_2-n_1}{r}.
    \]
    Then
    \[
      \frac{1.6}{p}-\frac{1.0}{5.0}=\frac{-0.6}{-3.0}=0.20.
    \]
    So \(1.6/p=0.40\), and \(p=\SI{4.0}{mm}\). The correct answer is
    \quizcorrectletter.
  \end{quizsolution}

  \begin{quizquestion}[
    id=optics-lens-010,
    topic=thin-lenses,
    difficulty=foundation,
    marks=1,
    correct=B,
    tags={converging-lens,focus},
    outcome={Identify converging lens focal sign}
  ]
    A converging thin lens has
    \choices
      {negative focal length and a virtual focus only.}
      {positive focal length and a real focus for incident parallel rays.}
      {zero focal length.}
      {positive focal length only for virtual images.}
      {negative focal length only when the object is inside \(F\).}
  \end{quizquestion}
  \begin{quizsolution}
    A converging lens brings parallel rays to a real focus, and the focal length
    is positive. The correct answer is \quizcorrectletter.
  \end{quizsolution}

  \begin{quizquestion}[
    id=optics-lens-011,
    topic=thin-lenses,
    difficulty=foundation,
    marks=1,
    correct=C,
    tags={diverging-lens,image-type},
    outcome={Recall the image formed by a diverging lens}
  ]
    For a real object in front of a diverging thin lens, the image is always
    \choices
      {real and inverted.}
      {real and upright.}
      {virtual and upright.}
      {virtual and inverted.}
      {formed at the far focal point.}
  \end{quizquestion}
  \begin{quizsolution}
    A diverging lens spreads rays so that their backward extensions meet on the
    object side. The image is virtual and upright. The correct answer is
    \quizcorrectletter.
  \end{quizsolution}

  \begin{quizquestion}[
    id=optics-lens-012,
    topic=thin-lenses,
    difficulty=applied,
    marks=1,
    correct=A,
    tags={thin-lens-equation,calculation},
    outcome={Use the thin lens equation}
  ]
    A converging lens has \(f=\SI{10}{cm}\). An object is
    \(\SI{15}{cm}\) from the lens. What is the image distance?
    \choices
      {\(+\SI{30}{cm}\)}
      {\(+\SI{6.0}{cm}\)}
      {\(-\SI{30}{cm}\)}
      {\(-\SI{6.0}{cm}\)}
      {\(+\SI{10}{cm}\)}
  \end{quizquestion}
  \begin{quizsolution}
    \[
      \frac{1}{i}=\frac{1}{10}-\frac{1}{15}=\frac{1}{30},
    \]
    so \(i=+\SI{30}{cm}\). The correct answer is \quizcorrectletter.
  \end{quizsolution}

  \begin{quizquestion}[
    id=optics-lensmaker-013,
    topic=thin-lenses,
    difficulty=applied,
    marks=1,
    correct=E,
    tags={lens-maker,calculation},
    outcome={Use the lens maker equation for a symmetric converging lens}
  ]
    A thin symmetric converging lens has \(n=1.65\) and \(f=\SI{4.0}{cm}\).
    If \(r_1=+r\) and \(r_2=-r\), what is \(r\)?
    \choices
      {\(\SI{1.3}{cm}\)}
      {\(\SI{2.6}{cm}\)}
      {\(\SI{4.0}{cm}\)}
      {\(\SI{4.6}{cm}\)}
      {\(\SI{5.2}{cm}\)}
  \end{quizquestion}
  \begin{quizsolution}
    Lens maker gives
    \[
      \frac{1}{f}=(n-1)\left(\frac{1}{r}-\frac{1}{-r}\right)
      =(n-1)\frac{2}{r}.
    \]
    Thus \(r=2(n-1)f=2(0.65)(4.0)=\SI{5.2}{cm}\). The correct answer is
    \quizcorrectletter.
  \end{quizsolution}

  \begin{quizquestion}[
    id=optics-twolens-014,
    topic=two-lens-systems,
    difficulty=challenge,
    marks=1,
    correct=B,
    tags={two-lens-system,calculation},
    outcome={Find the final image in a two-lens system}
  ]
    Lens 1 has \(f_1=+\SI{24}{cm}\), lens 2 has \(f_2=+\SI{9.0}{cm}\),
    and the lenses are \(\SI{10}{cm}\) apart. An object is
    \(\SI{6.0}{cm}\) to the left of lens 1. Where is the final image?
    \choices
      {\(\SI{8.0}{cm}\) to the left of lens 1}
      {\(\SI{18}{cm}\) to the right of lens 2}
      {\(\SI{18}{cm}\) to the left of lens 2}
      {\(\SI{9.0}{cm}\) to the right of lens 1}
      {At infinity}
  \end{quizquestion}
  \begin{quizsolution}
    Lens 1 gives \(i_1=-\SI{8.0}{cm}\). This point is
    \(\SI{18}{cm}\) to the left of lens 2, so \(p_2=\SI{18}{cm}\).
    Lens 2 gives
    \[
      \frac{1}{i_2}=\frac{1}{9.0}-\frac{1}{18}=\frac{1}{18}.
    \]
    Thus \(i_2=+\SI{18}{cm}\), to the right of lens 2. The correct answer
    is \quizcorrectletter.
  \end{quizsolution}

  \begin{quizquestion}[
    id=optics-twolens-015,
    topic=two-lens-systems,
    difficulty=challenge,
    marks=1,
    correct=D,
    tags={two-lens-system,magnification},
    outcome={Find the net magnification of a two-lens system}
  ]
    For the two-lens system with \(p_1=\SI{6.0}{cm}\),
    \(i_1=-\SI{8.0}{cm}\), \(p_2=\SI{18}{cm}\), and
    \(i_2=+\SI{18}{cm}\), what is the net magnification?
    \choices
      {\(+1\)}
      {\(-1\)}
      {\(+4/3\)}
      {\(-4/3\)}
      {\(-3/4\)}
  \end{quizquestion}
  \begin{quizsolution}
    \(m_1=-i_1/p_1=+8/6=4/3\), while
    \(m_2=-i_2/p_2=-18/18=-1\). Thus
    \(M=m_1m_2=-4/3\). The correct answer is \quizcorrectletter.
  \end{quizsolution}

  \begin{quizquestion}[
    id=optics-ray-016,
    topic=ray-diagrams,
    difficulty=foundation,
    marks=1,
    correct=A,
    tags={ray-diagram,thin-lens},
    outcome={Recall a principal ray for a converging lens}
  ]
    In a ray diagram for a converging lens, a ray initially parallel to the
    central axis emerges
    \choices
      {through the far focal point.}
      {through the near focal point before reaching the lens.}
      {without bending.}
      {back along itself.}
      {parallel to the central axis.}
  \end{quizquestion}
  \begin{quizsolution}
    A principal ray parallel to the central axis refracts through the far focal
    point of a converging lens. The correct answer is \quizcorrectletter.
  \end{quizsolution}

  \begin{quizquestion}[
    id=optics-magnifier-017,
    topic=optical-instruments,
    difficulty=applied,
    marks=1,
    correct=C,
    tags={simple-magnifier,angular-magnification},
    outcome={Compute simple magnifier angular magnification}
  ]
    A simple magnifier has \(f=\SI{5.0}{cm}\). Using
    \(m_u=\SI{25}{cm}/f\), what is its angular magnification?
    \choices
      {\(2\)}
      {\(4\)}
      {\(5\)}
      {\(20\)}
      {\(25\)}
  \end{quizquestion}
  \begin{quizsolution}
    \(m_u=25/5.0=5\). The correct answer is \quizcorrectletter.
  \end{quizsolution}

  \begin{quizquestion}[
    id=optics-microscope-018,
    topic=optical-instruments,
    difficulty=applied,
    marks=1,
    correct=E,
    tags={compound-microscope,magnification},
    outcome={Compute compound microscope magnification}
  ]
    A compound microscope has tube length \(s=\SI{16}{cm}\),
    \(f_{\text{ob}}=\SI{0.80}{cm}\), and
    \(f_{\text{ey}}=\SI{2.0}{cm}\). Estimate \(M\) using
    \(M\approx-s(\SI{25}{cm})/(f_{\text{ob}}f_{\text{ey}})\).
    \choices
      {\(-10\)}
      {\(-20\)}
      {\(-100\)}
      {\(-150\)}
      {\(-250\)}
  \end{quizquestion}
  \begin{quizsolution}
    \[
      M\approx-\frac{16(25)}{0.80(2.0)}
      =
      -\frac{400}{1.6}
      =
      -250.
    \]
    The correct answer is \quizcorrectletter.
  \end{quizsolution}

  \begin{quizquestion}[
    id=optics-telescope-019,
    topic=optical-instruments,
    difficulty=applied,
    marks=1,
    correct=B,
    tags={telescope,angular-magnification},
    outcome={Compute telescope angular magnification}
  ]
    A refracting telescope has \(f_{\text{ob}}=\SI{100}{cm}\) and
    \(f_{\text{ey}}=\SI{5.0}{cm}\). What is its angular magnification?
    \choices
      {\(+20\)}
      {\(-20\)}
      {\(+5\)}
      {\(-5\)}
      {\(-500\)}
  \end{quizquestion}
  \begin{quizsolution}
    \(m_u=-f_{\text{ob}}/f_{\text{ey}}=-100/5.0=-20\). The correct answer
    is \quizcorrectletter.
  \end{quizsolution}

  \begin{quizquestion}[
    id=optics-summary-020,
    topic=summary,
    difficulty=foundation,
    marks=1,
    correct=A,
    tags={conceptual-summary,signs},
    outcome={Choose the best general strategy for image problems}
  ]
    What is usually the best first move in an image-formation problem?
    \choices
      {Identify the optical element and predict whether the image should be real or virtual.}
      {Always assume \(i\) is positive.}
      {Always use the lens maker equation first.}
      {Ignore sign conventions until the final answer.}
      {Use telescope magnification for all lens problems.}
  \end{quizquestion}
  \begin{quizsolution}
    The physical setup should guide the sign choices before calculation. First
    identify the element and the expected image type. The correct answer is
    \quizcorrectletter.
  \end{quizsolution}
\end{quizbank}

\quizdefineversion{A}{%
  \quizselectall
  \quizshuffleoptions{34034}%
}

\quizuseversion{A}

\pagenumbering{gobble}
\makequiztitle
\pagenumbering{arabic}
\tableofcontents
\clearpage

\printquizquestions[2]

\begin{quizanswerkeycontent}
  \clearpage
  \printquizanswerkey
\end{quizanswerkeycontent}

\begin{quizsolutioncontent}
  \clearpage
  \printquizsolutions
\end{quizsolutioncontent}

\begin{quizteachercontent}
  \clearpage
  \printquiztopicreport
  \printquizteacherreport
\end{quizteachercontent}

\end{document}
